\begin{figure}[H]
\centering
\resizebox{\linewidth}{!}{%
\begin{tikzpicture}[
  card/.style={draw=npiNavy, thick, rounded corners=3pt, fill=white, minimum width=3.6cm, minimum height=1.8cm, align=center, font=\sffamily\scriptsize},
  arr/.style={->, thick}
]
  % A. Escalas de Weiser
  \node[font=\sffamily\bfseries\small\color{npiNavy}] at (5.2, 3.6) {A. Escalas espaciales de la computación ubicua (Weiser, 1991)};
  \node[card, draw=npiTeal, fill=npiTealLight] (tab) at (1.5, 2.0) {\textbf{\color{npiTeal}Tabs (cm)}\\[2pt]Etiquetas NFC, balizas BLE,\\pulseras y sensores\\llevables por usuario};
  \node[card, draw=npiCobalt, fill=npiCobaltLight] (pad) at (5.5, 2.0) {\textbf{\color{npiCobalt}Pads (dm)}\\[2pt]Smartphones y tabletas\\personales con contexto\\de identidad y postura};
  \node[card, draw=npiPlum, fill=npiPlumLight] (board) at (9.5, 2.0) {\textbf{\color{npiPlum}Boards (m)}\\[2pt]Pantallas públicas, mesas\\y puntos de información\\como el kiosco ETSIIT};

  \draw[<->, thick, color=npiSlate] (tab.east) -- (pad.west);
  \draw[<->, thick, color=npiSlate] (pad.east) -- (board.west);
  \node[font=\sffamily\scriptsize\color{npiSlate}] at (5.5, 0.55) {Escala creciente $\longrightarrow$ menor portabilidad, mayor visibilidad compartida};

  % Separador
  \draw[dashed, color=npiRule, thick] (11.75, 0.2) -- (11.75, 3.9);

  % B. Calm Technology: Centro vs Periferia
  \node[font=\sffamily\bfseries\small\color{npiTeal}] at (14.4, 3.6) {B. Atención (\textit{Calm Tech})};
  \draw[fill=npiPaper, draw=npiSlate, thick] (14.4, 1.8) circle (1.45cm);
  \draw[fill=npiAmberLight, draw=npiAmber, very thick] (14.4, 1.8) circle (0.68cm);
  \node[font=\sffamily\scriptsize\bfseries\color{npiAmber}] at (14.4, 1.8) {Centro};
  \node[font=\sffamily\tiny\color{npiSlate}, align=center] at (14.4, 2.8) {Periferia\\(vibración, luz)};
  \draw[<->, thick, color=npiTeal] (15.1, 1.8) -- (15.75, 1.8);
\end{tikzpicture}%
}
\caption{Taxonomía de escalas físicas de Mark Weiser y dinámica entre el centro y la periferia de la atención humana \parencite{weiser1991computer}.}
\label{fig:weiser_atencion}
\end{figure}
